%% RevTeX 4-2 is used as a neutral drafting format only.
%% Replace it with the target journal's class at submission.
\documentclass[aps,pra,reprint,superscriptaddress,amsmath,amssymb,floatfix]{revtex4-2}

\usepackage{graphicx}
\usepackage{xcolor}
\usepackage{amsthm}
\usepackage[expansion=false]{microtype}
\usepackage{hyperref}
\usepackage{orcidlink}

\hypersetup{
    colorlinks=true,
    linkcolor=blue,
    filecolor=magenta,
    urlcolor=blue,
    citecolor=blue,
    pdftitle={Computable Countermarkets and the Limits of Universal Trading},
    pdfauthor={Karl Svozil},
    pdfsubject={Computational, probabilistic, and financial limits of universal trading},
    pdfkeywords={trading, no-arbitrage, diagonalization, Halting Problem, rule inference, Martin-Lof randomness, algorithmic information theory, Busy Beaver, time reversal, market impact}
}

%% Mathematical Environments
\theoremstyle{plain}
\newtheorem{theorem}{Theorem}
\newtheorem{lemma}[theorem]{Lemma}
\newtheorem{corollary}[theorem]{Corollary}
\newtheorem{proposition}[theorem]{Proposition}
\newtheorem{observation}[theorem]{Observation}
\theoremstyle{definition}
\newtheorem{definition}{Definition}
\newtheorem{criterion}{Criterion}
\theoremstyle{remark}
\newtheorem{remark}{Remark}

\begin{document}

\title{Computable Countermarkets and the Limits of Universal Trading}

\author{Karl Svozil\,\orcidlink{0000-0001-6554-2802}}
\affiliation{Institute for Theoretical Physics, TU Wien, Wiedner Hauptstra{\ss}e 8-10/136, A-1040 Vienna, Austria}
\email{karl.svozil@tuwien.ac.at}
\homepage{http://tph.tuwien.ac.at/~svozil}

\date{September 19, 2026}

\begin{abstract}
Every universal-trading claim pairs a trader with a market---a path, generator,
or law. For any total deterministic computable trader, its code yields a fixed
positive rational countermarket for one discounted asset: every proportional
return is either \(+\epsilon\) or \(-\epsilon\), every active position loses,
and inactivity earns zero. Thus the obstruction persists when flat paths are
excluded. The model is frictionless and self-financing, with rational orders
and no computation deadline. Gold-style learning cannot
identify every computable binary market rule from history. Separately,
Turing-universal generators make certification of unbounded-future events
undecidable; Busy-Beaver growth defeats every computable description-size
waiting schedule. A passive Martin-L\"of-random record is incompressible and
prevents effective test-capital processes from becoming unbounded;
no-arbitrage supplies a distinct financial boundary. Together these results
form an expository taxonomy: repeatable success requires a market restriction,
benchmark, risk premium, or informational advantage, while time reversal and
within-class permutations supply simple stress tests.
\end{abstract}

\maketitle

%%------------------------------------------------------------------------------
\section{Introduction}
\label{sec:intro}
%%------------------------------------------------------------------------------

Throughout this paper the elementary configuration has two sides: a
\emph{trader} and a \emph{market}. At each date the trader converts observed
history into a position; the market supplies the next price. Depending on the
question, the market side is a realized path, a path-generating program, a
probability law, or a price process constrained by no-arbitrage. Holding the
trader fixed while changing the admitted market side reveals which universal
claim is at issue.

A strategy is \emph{universal} here only if, after discounting by a specified
numeraire, it produces strictly positive terminal gain from zero initial
capital on every admitted trajectory at a specified finite horizon. This
differs from nonnegative gain everywhere with positive gain somewhere, eventual
profit at a path-dependent date, expected profit under a law, and small regret
to a benchmark. Zero cost, self-financing,
admissibility, and discounted gain distinguish this claim from positive wealth
obtained by investing positive capital. Its quantifier order is
\begin{equation}
  \exists\,\sigma\ \forall\,\mathcal E:
  \qquad \sigma\ \hbox{wins against }\mathcal E,
  \label{eq:universalquantifiers}
\end{equation}
whereas diagonalization establishes
\begin{equation}
  \forall\,\sigma\ \exists\,\mathcal E_\sigma:
  \qquad \sigma\ \hbox{does not win against }\mathcal E_\sigma.
  \label{eq:counterquantifiers}
\end{equation}

For a fixed total deterministic computable trader, its code can be embedded in
a total program that recursively chooses the next bounded proportional price
move opposite to each nonzero position and a nonzero move during inactivity.
The resulting countermarket is one
fixed computable path; any finite prefix can be generated offline. A trader that
never takes a position earns zero, which already refutes a guarantee of strict
profit. This is the trading counterpart of the next-bit construction
\(x_0=1-F(\varnothing)\) and
\(x_t=1-F(x_0,\ldots,x_{t-1})\) for \(t\geq1\): it defeats the fixed proposed
rule, not every rule on one exceptional sequence.

Three complementary limits use different market information. Gold's learner
sees an accumulating history and must eventually stabilize on a correct
program; no effective learner identifies every binary total recursive time
function, although restricted classes such as the primitive-recursive time
functions are identifiable in the limit~\cite{go-65,go-67}. A source-code
procedure that settled every future event of Turing-universal generators would
instead decide the Halting Problem, while Busy-Beaver growth excludes a
computable worst-case waiting horizon. Finally, relative to a specified
computable measure, a Martin-L\"of-random record escapes every effective null
test, has incompressible prefixes, and prevents effective nonnegative
test-capital from becoming unbounded. No-arbitrage supplies a separate
pricing-measure boundary.

Repeatable practical success must therefore identify structure absent from the
unrestricted market class: persistence, mean reversion, a risk premium, a
benchmark, lawful information, or another distributional asymmetry. Large,
crowded, or public strategies may also change the distribution through price
impact and adaptation. This falsifiable restriction on the market side is the
operational meaning of ``insight'' used below.

\begin{table*}[t]
\caption{\label{tab:claimmap}Distinct claims organized by the trader--market
configuration. They are complementary, not interchangeable.}
\centering
\begin{ruledtabular}
\begin{tabular}{p{0.18\textwidth}p{0.27\textwidth}p{0.24\textwidth}p{0.23\textwidth}}
\textbf{Mechanism} & \textbf{Market side} & \textbf{Formal scope}
& \textbf{Conclusion} \\
\colrule
Diagonal\newline\mbox{countermarket}
& Fixed total program constructed from the fixed trader's code
& Pathwise worst case:
\(\forall\sigma\,\exists P_\sigma\)
& No total deterministic computable trader wins on every computable price path \\
Gold identification\newline in the limit
& Accumulating computable history, or interaction with a black box
& Success is eventual stabilization on a correct program, without a known
stopping date
& No effective learner identifies every computable binary time function;
restricted classes can be identified \\
Halting reduction and\newline Busy Beaver
& Total program or generator, under an explicit promise
& Individual encoded computation over an unbounded horizon
& No general computable profit or market-event decider, or computable
description-size waiting horizon \\
Algorithmic\newline randomness
& Passive binary record under the computable fair-coin reference law
& Each Martin-L\"of-random path; such paths have measure one
& No constructive supermartingale, hence no computable fair-game martingale,
becomes unbounded \\
No-arbitrage
& Stochastic price process admitting a coherent pricing measure
& Almost-sure payoff conditions under an ELMM
& No admissible zero-cost classical arbitrage \\
Cover universality
& Passive bounded sequence and a stated benchmark class
& Pathwise, but benchmark-relative and asymptotic
& Vanishing per-period regret, not guaranteed positive profit \\
\end{tabular}
\end{ruledtabular}
\end{table*}

The explicit diagonal countermarket is the constructive core of the paper.
Table~\ref{tab:claimmap} organizes the complementary limits into a taxonomy that
keeps distinct several statements commonly summarized as ``no universal
strategy.'' The construction and its proofs are expository reformulations
rather than a claim of a newly discovered impossibility theorem.
The analysis connects this quantifier map to the literature on unpredictable
prices~\cite{Bachelier1900,Samuelson1965,Fama1970}, Gold learning, computability
and randomness, no-arbitrage and NFL, time-reversal and permutation stress
testing, the Wheel, Cover universality, and conditional FAPP strategies.

%%------------------------------------------------------------------------------
\section{The Trader and the Market: Computability and Randomness}
\label{sec:computability}
%%------------------------------------------------------------------------------

We now formalize the trader--market pair. The trader is a program
\(M_\sigma\) mapping each observed finite price history to a position; the
market side supplies the price history. Four specifications lead to four
different boundaries: a fixed computable countermarket tailored to a trader
defeats its claim of pathwise universality; Gold's history protocol precludes
identifying every computable rule in the limit; a computationally universal
class of market generators makes uniform certification of unbounded future
events undecidable; and a passive algorithmically random record defeats
effective fair betting.

\subsection{A Computable Countermarket for Every Trader}

Throughout the discrete constructions, \(P_t=S_t/B_t\) is the discounted
price of one asset, where \(S_t\) is its price in a fixed currency and
\(B_t>0\) is the value of one unit of the chosen risk-free bank account.
For example, if both initially equal \(100\), and later \(S_1=121\) and
\(B_1=110\), the discounted price rises from \(1\) to \(1.1\).
This measures performance relative to the bank account, not inflation.
The position \(w_t\) is a number of shares chosen before \(P_{t+1}\) is
observed. Trading is self-financing: changes in holdings are funded through
the existing share and bank balances, with no external cash injections.
More explicitly, holding \(w_t\) shares and the residual
\(V_t-w_tP_t\) in the discounted bank account gives
\[
  V_{t+1}=V_t+w_t(P_{t+1}-P_t).
\]
We omit transaction costs, borrowing limits, leverage caps, and solvency
constraints. The counterexample therefore also applies to any subclass of
these strategies satisfying additional restrictions.

Let
\(M_\sigma\colon\mathbb Q_{>0}^{<\omega}\to\mathbb Q\)
be a total computable function on finite rational histories, encoded in the
usual effective way, with \(w_t=M_\sigma(P_0,\ldots,P_t)\).
Totality requires an
action in finite computation time on every input; it supplies no prescribed
running-time bound or guarantee of meeting a trading deadline. Rational outputs
model finite-precision orders and, crucially, make comparison with zero
decidable. Let \(\mathcal P_{\mathrm{comp}}^{+}\) denote the positive
rational-valued computable price paths.

\begin{samepage}
\begin{theorem}[Computable Diagonal Countermarket]
  \label{thm:counterpath}
  Fix rational \(P_0>0\) and \(0<\epsilon<1\). For every total deterministic
  computable strategy \(M_\sigma\), there exists a fixed path
  \(P^{M_\sigma}\in\mathcal P_{\mathrm{comp}}^{+}\), starting at \(P_0\), with
  \(P_{t+1}/P_t\in\{1-\epsilon,1+\epsilon\}\) at every date, on which its
  self-financing gain
  \[
    G_T=\sum_{t=0}^{T-1}w_t(P_{t+1}-P_t)
  \]
  is non-positive at every finite horizon \(T\). It is strictly negative after
  the strategy has taken a nonzero position at least once.
\end{theorem}
\end{samepage}

\begin{proof}
Choose rational \(P_0>0\) and rational \(0<\epsilon<1\). Recursively compute
\(w_t=M_\sigma(P_0,\ldots,P_t)\) and set
\begin{equation}
P_{t+1}=
\begin{cases}
P_t(1-\epsilon),&w_t>0,\\
P_t(1+\epsilon),&w_t<0,\\
P_t(1+\epsilon(-1)^t),&w_t=0.
\end{cases}
\label{eq:adversary}
\end{equation}
Because \(M_\sigma\) is total and rational-valued, this recursion is a total
program for one fixed infinite price path. Its prices remain positive and
computable, and every return has magnitude \(\epsilon\). In each period,
\(w_t(P_{t+1}-P_t)\leq0\), with strict inequality when \(w_t\neq0\).
Summing proves the claim at every finite horizon.
\end{proof}

For a concrete example, fix the rule that holds one share for the first
period, shorts one share for the second, and then holds no shares. With
\(P_0=100\) and \(\epsilon=1/10\), Eq.~\eqref{eq:adversary} produces
\begin{align*}
  (w_0,w_1,w_2)&=(1,-1,0),\\
  (P_0,P_1,P_2,P_3)&=(100,90,99,108.9).
\end{align*}
The successive discounted gains are \(-10,-9,0\), so \(G_3=-19\), all in
bank-account units. The prices are constructed after fixing the rule; they
are illustrative, not market data.

\begin{remark}[Logical and Economic Scope of the Countermarket]
\label{rem:counterpath-scope}
Theorem~\ref{thm:counterpath} uses the trader's \emph{code}, not a physical
market's observation of a deployed order. Once \(M_\sigma\) is fixed, the
countermarket program is fixed and can generate any desired finite prefix
offline. It is therefore a lawful, deterministic, computable counterexample,
although not necessarily an efficient one computationally. If the market side
is regarded as an opponent whose objective is to prevent this trader's gain,
the rule is rational relative to that objective. It is not thereby shown to be
an economic equilibrium or an arbitrage-free price process, and another trader
may profit on it. Such a trader has, in turn, its own computable countermarket.
\end{remark}

This is a direct strategy-relative diagonal construction. The proposed
universal trader is used to compute an object outside its success set: long is
paired with down, short with up, and neutrality with an alternating nonzero
return. Just as
the complement sequence defeats a purported predictor of every computable
sequence, \(P^{M_\sigma}\) defeats the purported trader on every horizon.
Enumeration of all programs is unnecessary because the universal claim itself
supplies the program to be diagonalized against. The quantifiers are
\(\forall M_\sigma\,\exists P^{M_\sigma}\). A constant discounted price path
already gives every self-financing rule zero trading gain and refutes strict
profit on a class containing that path. The stronger content of
Theorem~\ref{thm:counterpath} is that every nonzero position of the selected
trader incurs a strict loss even on the class with fixed return magnitude
\(\epsilon\), which contains no flat path. It does not supply a path on which every strategy
strictly loses: permanently long and permanently short positions cannot both
lose on the same one-step price move.

\subsection{Learning a Market Rule from History: Gold's Limit}

A trader normally receives neither the market generator's source code nor a
proof of its law. The trader sees an expanding record
\(x_0,x_1,\ldots\) and must infer from each finite prefix how it will continue.
Gold's function-identification model formalizes exactly this information flow.
At time \(n\), an effective learner \(L\) maps
\((x_0,\ldots,x_n)\) to the index of a Turing machine. It
\emph{identifies \(x\) in the limit} if, from some finite date onward, every
guess is the same index and that machine computes the whole sequence. The
learner need not know when its last change of mind has occurred
\cite{go-65,go-67}.

Let \(\mathrm{REC}_2\) denote the total computable binary sequences. Gold's
negative theorem for binary-valued time functions has the trader--market
quantifier order, with \(L\) ranging over effective learners,
\begin{equation}
  \begin{aligned}
  &\forall L\ \exists x\in\mathrm{REC}_2: \\
  &\quad L(x_0,\ldots,x_n)\text{ does not stabilize}\\
  &\quad \text{on a program for }x .
  \end{aligned}
  \label{eq:goldquantifiers}
\end{equation}
The adverse history depends on the learner. This is closely analogous to
\(\forall M_\sigma\,\exists P^{M_\sigma}\), although Gold's failure criterion
is failure to converge to a fixed correct program name, not financial loss.

The corresponding claim about literal next-move forecasting is sharper and
has the short direct diagonal proof already used above. If a total computable
predictor \(F\) outputs one bit from every finite binary history, the recursive
definition
\[
  x_0=1-F(\varnothing),
  \qquad
  x_t=1-F(x_0,\ldots,x_{t-1})\quad(t\geq1)
\]
produces a computable market record on which \(F\) is wrong at every date.
Thus the proposed statement---every total next-move predictor has a computable
countermarket---is correct, but it is this direct prediction diagonalization,
not the literal statement of Gold's identification theorem. Nonconvergence of
program names alone does not count the learner's next-step errors.

The opposite-bit construction also appears explicitly in Legg's analysis of
universal prediction~\cite{Legg2006}. The qualifier \emph{computable} is
essential. Solomonoff induction defines an uncomputable universal predictive
ideal; for computable probabilistic environments its conditional predictions
converge, with quantitative error bounds developed by Hutter
\cite{solomonoff,solomonoff-II,solomonoff1,Hutter2001}. It therefore evades,
rather than contradicts, the diagonal result: it is not a total executable
trading rule, and predictive convergence is not Gold's requirement of
stabilizing on one exact generating program. For probability-valued forecasts,
V'yugin gives a further, distinct impossibility: a randomized construction can
make any partial computable forecaster fail calibration with probability
arbitrarily close to one~\cite{Vyugin2006}. Calibration, next-bit error, and
identification in the limit are three different criteria.

Gold's results also show why the restriction on the market side matters. Every
uniformly effective indexed family of total computable functions is
identifiable in the limit: a learner selects the least index consistent with
the observed prefix, and eventually eliminates every smaller incorrect index.
In particular, primitive-recursive time functions form such a family. In Gold's interactive
black-box model, where the eventual hypothesis must reproduce future
input--output behavior, finite-state boxes are weakly identifiable in the
limit, whereas the class of primitive-recursive black boxes is not
\cite{go-65}. These are different object classes, not contradictory claims.
Gold's 1967 language model further makes explicit that learnability depends on
the admitted class, the presentation of the evidence, and the allowed names
\cite{go-67}. For trading, a forecast can therefore be justified only relative
to a specified restriction on possible markets.

\subsection{Universal Computation and the Halting Barrier}

Gold's learner works from observations and may revise its hypothesis
indefinitely. The next results ask a separate, stronger operational question:
can one terminating program inspect arbitrary encoded traders or market
generators and correctly decide their unbounded future behavior? The preceding
diagonal construction already disproves a trader that wins on every computable
path and does not invoke the Halting Problem. The source-code
event-certification question does.

\begin{theorem}[Undecidability of Eventual Profit]
  \label{thm:profit-undecidable}
  Fix the positive computable price path \(P_t=t+1\), \(t\in\mathbb{N}\).
  There is no algorithm which, under the promise that its input code computes
  a total trading program \(M_\sigma\), correctly decides for every such code
  whether
  \[
    \exists T\geq1:\qquad
    G_T(M_\sigma;P)=
    \sum_{t=0}^{T-1}w_t(P_{t+1}-P_t)>0.
  \]
\end{theorem}

\begin{proof}
Assume such a profit certifier exists. Given an arbitrary Turing machine
\(A\) and input \(x\), construct a trading program \(M_{A,x}\) as follows.
At date \(t\), simulate \(A(x)\) for \(t+1\) steps. If it has halted within
those steps, set \(w_t=1\); otherwise set \(w_t=0\). Every date requires only
a finite simulation, so \(M_{A,x}\) is total. Along \(P_t=t+1\), every long
position earns one unit in the next period. Hence \(G_T>0\) for some \(T\) if
and only if \(A(x)\) eventually halts. A certifier for eventual trading profit
would therefore decide the Halting Problem, which is impossible
\cite{turing-36}.
\end{proof}

The reduction can also be placed directly on the market side, which is the
form most closely connected to source-code event certification.

\begin{theorem}[Undecidability of a Market Event]
  \label{thm:market-undecidable}
  There is no algorithm which, under the promise that its input code computes
  a total positive-price generator \(G\), correctly decides for every such
  code whether its generated path \(P^G\) ever satisfies \(P_t^G>P_0^G\).
\end{theorem}

\begin{proof}
Given a Turing machine \(A\) and input \(x\), define the total generator by
\(P_0^{A,x}=1\) and, for \(t\geq1\),
\begin{equation}
  P_t^{A,x}
  =1+\mathbf{1}\{A(x)\text{ halts within }t\text{ steps}\},
  \qquad t\geq1 .
  \label{eq:market-halting}
\end{equation}
Each quoted price requires only a finite \(t\)-step simulation, so the
generator is total and computable. Its initial price is one, and it ever quotes
a higher price if and only if \(A(x)\) halts. A universal decider for the stated
market event would therefore decide the Halting Problem.
\end{proof}

Theorem~\ref{thm:market-undecidable} makes ``universal computation'' precise:
the admissible generators can emulate an arbitrary machine and expose its
outcome through a price event. No algorithm can then certify every future
event of every such market. Both reductions concern promise domains: the
putative decider must halt and be correct on every promised total input.
Positive instances are semidecidable by simulation until the event occurs;
there is no terminating general procedure for certifying its perpetual absence.

This is not Gold's theorem. The reduction receives source code and must halt
with a yes-or-no answer; Gold's learner receives accumulating data and need
only stabilize eventually. Indeed, the paths in
Eq.~\eqref{eq:market-halting} are identifiable in the limit: guess the flat
path, then switch permanently if the price becomes two. No finite date,
however, certifies that the initial guess is final. Thus halting excludes
uniform certification, whereas Gold excludes identification in the limit for
the class of all binary total recursive time functions. Restricted model
classes may still be learnable
\cite{blum75blum,angluin:83,ad-91,2016-AlexanderKarlSvozil-ml,moore}.

\paragraph{Busy-Beaver horizons: finite glimpses and noncomputable waiting}

Fix a step-counted universal machine \(V\) with an effective compiler that adds
only fixed overhead for a given program schema, and define the runtime Busy
Beaver function by
\begin{equation}
  \operatorname{BB}(n)
  =
  \max\{\operatorname{time}_V(p):
  |p|\leq n\text{ and }V(p)\text{ halts}\},
  \label{eq:busybeaver}
\end{equation}
with value zero if the set is empty. The input \(n\) is a program-description
budget, not a date; \(\operatorname{BB}(n)\) is a runtime, which becomes a
possible trading date when the computation is embedded in a market generator.
A small increase in rule length may therefore require waiting beyond every
computably prescribed scale.

Each \(\operatorname{BB}(n)\) is finite, and selected low values can sometimes
be proved by exhaustive simulation plus proofs that the remaining programs do
not halt. This is not uniform ``computability for small arguments'': any finite
table can be hard-coded, and its apparent pattern licenses no extrapolation. If
a total computable \(h\) upper-bounded \(\operatorname{BB}\), simulating every
program of length at most \(n\) for \(h(n)\) steps would decide halting.
Eventual domination uses a further program-schema argument. For each \(n\), a
program of \(O(\log n)\) bits can encode \(n\), compute \(h(n)\), execute more
than \(h(n)\) dummy steps, and halt. For all sufficiently large \(n\), its
compiled \(V\)-code has length at most \(n\), so
\(\operatorname{BB}(n)>h(n)\). Thus, up to the fixed machine convention,
\(\operatorname{BB}\) eventually exceeds every total computable
function~\cite{rado,chaitin-bb}.

Calude, Dinneen, and Shu illustrate the finite/infinite divide. For one explicit
universal self-delimiting register machine, computation and mathematical proof
settled halting through 84-bit programs and certified 64 initial bits of its
halting probability \(\Omega_U\). They did not compute Busy Beaver values, and
the work needed to recover the relevant halting programs from longer supplied
prefixes of \(\Omega_U\) outruns every computable bound. Their result is a
formidable finite glimpse, not a scalable computation of an uncomputable object
\cite{2002-glimpseofran}.

Applied to Eq.~\eqref{eq:market-halting}, a short generator can quote
\(P_t=1\) for a Busy-Beaver-scale interval and then quote \(P_t=2\). Across
encoded programs of length at most \(n\), the last halt can occur at
\(\operatorname{BB}(n)\), up to fixed wrapper overhead. Knowing that horizon
would make continued silence a certificate of no later move. Equivalently,
with \(\mathrm{TotComp}\) denoting total computable size-dependent schedules,
\begin{equation}
  \begin{aligned}
  \forall h\in\mathrm{TotComp}\ \exists p:\quad
    &V(p)\text{ halts},\\
    &\operatorname{time}_V(p)>h(|p|).
  \end{aligned}
  \label{eq:bb-horizon}
\end{equation}
If \(c\) is the wrapper's fixed description-length overhead, it is absorbed by
the computable envelope \(H(n)=\max_{k\leq n+c}h(k)\). Every quote at a
specified finite date remains
computable by bounded simulation; the generator neither computes Busy Beaver
nor selects its maximizer uniformly. Noncomputability lies in certification
across the family. Each realized path is deterministic and computable, hence
not Martin-L\"of random under fair-coin coding.

This barrier is stronger than sensitive dependence or visual complexity alone
in a precise decision-theoretic sense. A finite contractive iterated-function
system or chaotic map with computable data can generate an intricate phenotype
from a short rule, yet Turing universality additionally creates questions about
the unbounded future for which no terminating general algorithm exists. The
categories can overlap when a dynamical system performs universal computation.
This is a worst-case claim across generator descriptions, not evidence that
typical financial markets exhibit Busy-Beaver delays
\cite{barnsley:88,moore,calude:037103}.

These obstructions concern unbounded horizons. At any specified finite horizon,
total programs can still be run and their realized prices and gains computed.

\subsection{Worst-Case Value and Private Randomization}

Returning to the direct countermarket, its finite-horizon value can be stated
explicitly, and the role of a trader's hidden randomization can be separated
from deterministic pathwise opposition.

\begin{corollary}[Maximin Value]
  \label{cor:maximin}
  At every finite horizon \(T\), if the neutral strategy is allowed, then
  \[
    \sup_{\text{\rm total computable }M_\sigma}
    \inf_{P\in\mathcal P_{\mathrm{comp}}^{+}}
    G_T(M_\sigma,P)=0.
  \]
\end{corollary}

\begin{proof}
The neutral strategy guarantees zero, while
Theorem~\ref{thm:counterpath} gives every strategy a computable countermarket
with gain at most zero.
\end{proof}

\begin{proposition}[Private Randomization Does Not Create a Guarantee]
  \label{prop:randomized}
  Fix a finite horizon \(T\), a rational \(P_0>0\), and a rational
  \(\epsilon\in(0,1)\). There is a single passive computable market law,
  independent of the trader, under which every predictable randomized strategy
  with integrable gains has expected gain zero. Specifically, let
  \[
    P_{t+1}=P_t(1+\epsilon\xi_{t+1}),
    \qquad t=0,\ldots,T-1,
  \]
  where \((\xi_t)_{t\geq1}\) are independent fair signs, also independent of
  the trader's private random seed. If \(w_t\) is measurable with respect to
  the information available when the position is chosen and
  \(w_t(P_{t+1}-P_t)\) is integrable, then
  \[
    G_n=\sum_{t=0}^{n-1}w_t(P_{t+1}-P_t),
    \qquad n=0,\ldots,T,
  \]
  is a martingale. In particular, \(\mathbb E[G_T]=0\).
\end{proposition}

\begin{proof}
Let \(R\) denote the trader's private seed and set
\(\mathcal G_t=\sigma(P_0,\ldots,P_t,R)\). Allowing \(\mathcal G_t\) to contain
all of \(R\) only strengthens the claim. Predictability makes \(w_t\)
\(\mathcal G_t\)-measurable, while \(\xi_{t+1}\) is independent of
\(\mathcal G_t\) and has mean zero. Hence
\[
 \mathbb E\!\left[
   w_t(P_{t+1}-P_t)\mid\mathcal G_t
 \right]
 =\epsilon P_t w_t\,
   \mathbb E[\xi_{t+1}\mid\mathcal G_t]
 =0.
\]
The integrability assumption therefore makes \((G_n)_{n=0}^T\) a martingale,
and \(\mathbb E[G_T]=\mathbb E[G_0]=0\).
\end{proof}

Private randomization prevents a deterministic countermarket based only on the
public history from mechanically opposing the unseen current draw, but it does
not yield strictly positive expected gain in every environment. More precisely,
under the proposition's integrability assumptions it excludes
\(\forall P:\mathbb P_R(G_T(P,R)>0)=1\) at the fixed horizon \(T\).
Each of the finitely many fair-sign prefixes has positive probability and a
computable infinite continuation; such a guarantee would imply
\(G_T>0\) almost surely jointly, hence \(\mathbb E[G_T]>0\), a contradiction.
It does not exclude a high win probability, profit under another law, or
benchmark-relative success.

\paragraph{Optional online-game reading}

The recursion in Theorem~\ref{thm:counterpath} also defines a finite
perfect-information game: the trader announces \(w_t\), the environment chooses
the next return, and Corollary~\ref{cor:maximin} gives the trader-first maximin
value zero. This reading is optional because the fixed deterministic trader's
countermarket program already generates the same path offline. Borel
determinacy concerns infinite games with Borel payoff sets and does not by
itself identify a winner here~\cite{Martin1975}. Hidden private randomization is
a different information structure: Proposition~\ref{prop:randomized} proves
zero expected gain under one passive computable fair-sign law, not a mixed-game
value without specified strategy spaces.

\subsection{A Passive Market: Martin-L\"of Randomness}

\paragraph{When does process randomness yield a random record?}

Process randomness describes how individual bits are generated and depend on
their past. Martin-L\"of (ML) randomness is a property of the resulting
\emph{infinite sequence} relative to a specified probability law.
\emph{Independent fair bits produce a fair-coin ML-random sequence with
probability one.} Equivalently, the sufficient process condition is
\begin{equation}
  \mathbb P(X_{n+1}=1\mid X_1,\ldots,X_n)=\tfrac12
  \quad\text{a.s. for every }n\geq0,
  \label{eq:conditional-fairness}
\end{equation}
including the first bit at \(n=0\).

Independence alone is insufficient for \emph{fair-coin} randomness: independent
bits of fixed bias \(p\neq1/2\) almost surely have limiting frequency \(p\).
Fair one-bit probabilities alone are also insufficient: repeating one fair
bit forever gives fair marginals but a constant, predictable record.
Independence is not necessary for ML randomness relative to the process's
\emph{own} computable law, which may include bias or dependence.

For FAPP use, the conclusion is therefore conditional on a justified source
model, for example independent fair trials. Tests for bias and serial
dependence can challenge that model over observed scales; passing them cannot
certify independence indefinitely or the ML randomness of an infinite output.
An almost-sure conclusion allows exceptional records of probability zero.
A fixed deterministic pseudorandom generator with a finite seed cannot produce
a fair-coin ML-random infinite sequence, however well its finite outputs test.

Here is the precise classical result. A computable law \(\mu\) has cylinder
probabilities \(\mu([s])\) uniformly approximable to any requested accuracy,
where \([s]\) comprises sequences beginning with \(s\). An ML test is a
uniformly computably enumerable sequence of open sets \((U_k)\) with
\(\mu(U_k)\leq2^{-k}\). A record is \(\mu\)-ML-random if it avoids
\(\bigcap_kU_k\) for every such test.

\begin{proposition}[Classical Measure-One Result, Martin-L\"of]
  \label{prop:process-ml}
  If a binary process has computable joint law \(\mu\), then
  \begin{equation}
    \mathbb P\bigl((X_1,X_2,\ldots)\in\mathrm{MLR}_\mu\bigr)=1.
    \label{eq:process-ml}
  \end{equation}
  This includes independent fair trials, computable Bernoulli biases, and
  computable dependent laws~\cite[Sec.~IV]{MartinLof1966}.
\end{proposition}

\begin{proof}
Each test's failure set has measure zero. There are countably many effective
tests, so their union still has measure zero.
\end{proof}

The law must be specified: randomness relative to an atomic law need not mean
fresh uncertainty at each trial. For example, either constant record is random
for the law that repeats one fair bit forever.

\paragraph{Borel normality is weaker}
\label{par:normality}

A sequence is Borel normal in base \(b\) when every word of length \(k\)
has limiting frequency \(b^{-k}\), for every \(k\geq1\).
ML randomness for uniform base-\(b\) digits implies normality, but the converse
fails. Champernowne's decimal
\(0.123456789101112\ldots\), formed by concatenating the positive integers,
is normal in base ten yet computable, hence not ML-random for the uniform
decimal law~\cite{ch1}. Normality is difficult to establish for a given
infinite record: one must control all block lengths in the infinite limit.
No finite frequency check proves it, since any finite prefix admits both
normal and nonnormal continuations. A proof from a construction, as for
Champernowne, can establish normality; extensive sampling alone cannot.

\paragraph{Compression and effective betting}

For a universal prefix-free machine \(U\), let
\begin{equation}
  K_U(s)=\min\{|p|:U(p)=s\}.
  \label{eq:prefixcomplexity}
\end{equation}
The classical Levin--Schnorr characterization says that a binary record \(x\)
is fair-coin ML-random exactly when, for some constant \(c\),
\begin{equation}
  K_U(x{\upharpoonright}n)\geq n-c\quad\text{for every }n.
  \label{eq:levinschnorr}
\end{equation}
Its prefixes cannot be compressed by unbounded savings~\cite{Chaitin-1974,calude:02,Nies2009}.
The complexity \(K_U\) is uncomputable; unsuccessful finite compression tests
do not certify this condition.

A nonnegative capital function \(d\) is a fair-coin supermartingale if
\begin{equation}
  d(s)\geq\frac{d(s0)+d(s1)}2;
  \label{eq:computablemartingale}
\end{equation}
equality gives a martingale. Lower semicomputability means that its values
are uniformly approximable from below by increasing computable rationals.

\begin{theorem}[Classical Ville Inequality]
  \label{thm:ville}
  For a nonnegative supermartingale with initial capital \(d(\varnothing)\),
  under fair-coin law \(\lambda\),
  \begin{equation}
    \lambda\{x:\sup_n d(x{\upharpoonright}n)\geq a\}
    \leq d(\varnothing)/a\qquad(a>0).
    \label{eq:ville}
  \end{equation}
\end{theorem}

Stopping at the first crossing and taking limits proves the inequality
\cite{Ville1939}. Effectiveness turns its null exceptional sets into ML tests.

\begin{theorem}[Classical Effective Martingale Characterization]
  \label{thm:effective}
  A binary record is fair-coin ML-random if and only if every
  lower-semicomputable nonnegative supermartingale has bounded capital on it.
  In particular, every total computable nonnegative fair betting strategy
  remains bounded on each such record.
\end{theorem}

\begin{proof}
For a fixed finite-valued supermartingale \(d\), choose an integer \(c\geq0\)
with \(d(\varnothing)\leq2^c\), and put
\[
  U_k=\{x:\exists n,\ d(x{\upharpoonright}n)>2^{k+c}\}.
\]
These sets are uniformly effectively open and have measure at most \(2^{-k}\)
by Ville's inequality. The integer \(c\) is chosen existentially for this
\(d\); no uniform upper-bound algorithm is asserted. Thus unbounded capital
would fail an ML test. Conversely, each ML test admits a lower-semicomputable
supermartingale unbounded on its failure set~\cite{calude:02,Nies2009}.
\end{proof}

For a computable law \(\mu\), the corresponding condition is
\(d(s)\mu([s])\geq d(s0)\mu([s0])+d(s1)\mu([s1])\).
Such a test process need not supply an executable wager by a trading deadline.
Nor is it automatically a self-financing portfolio: that interpretation
requires suitable traded payoffs and capital constraints. A computable law
alone is not an equivalent local martingale measure (ELMM); an ELMM must also
be equivalent to the physical law and make discounted traded prices local
martingales.

\begin{corollary}[Asymptotic Accuracy of Computable Prediction]
  \label{cor:prediction-edge}
  Let \(x\) be fair-coin ML-random and
  \(F\colon\{0,1\}^{<\mathbb N}\to\{0,1\}\) a total computable predictor.
  Then
  \[
    A_T=\frac1T\sum_{t=1}^T\mathbf1\{F(x_1\cdots x_{t-1})=x_t\}
    \longrightarrow\frac12.
  \]
\end{corollary}

\begin{proof}
Set \(e_t=\mathbf1\{F(x_1\cdots x_{t-1})\neq x_t\}\). The map
\(x\mapsto e\) is a total computable, fair-coin-measure-preserving bijection:
each error prefix reconstructs exactly one input prefix recursively.
Computable measure-preserving maps preserve ML randomness~\cite{Rute2016}.
Thus \(e\) is ML-random and its frequency of ones tends to \(1/2\), giving
the asserted accuracy limit.
\end{proof}

These statements concern an infinite record and nonnegative test capital.
Proposition~\ref{prop:randomized} instead concerns integrable gains at a fixed
finite horizon. Neither excludes finite winning streaks, and direction-only
coding omits return magnitudes, timing, and costs relevant to trading profit.

\paragraph{Recent extensions}

Huang, J.~H. Lutz, N. Lutz, and Migunov (2026) extend the martingale,
effective-measure, and complexity characterizations to individual
continuous-time Markov trajectories. For chemical reaction networks, random
trajectories with bounded molecular counts cannot have infinitely many
reactions in finite time~\cite{HuangLutzLutzMigunov2026}.
Huttegger, Walsh, and Zaffora Blando (2026) characterize ML randomness through
summable one-step conditional relative-entropy discrepancies against computable
comparison laws with bounded prefix relative entropy, assuming computable
full-support laws~\cite{HutteggerWalshZafforaBlando2026}.
These extend the analysis of dependent processes and forecast agreement;
the elementary almost-sure bridge above already belongs to Martin-L\"of's
1966 framework.

\subsection{Ramsey Regularity and the Finite-Pattern Trap}
\label{sec:ramseytrap}

A large market record can be pattern-rich even when it contains no exploitable
law. Ramsey theory makes this finite combinatorial point without assuming a
probability model. For fixed positive integers \(c,m,q\), there is a finite
number \(R_q(m;c)\) such that every \(c\)-coloring of the \(q\)-element subsets
of any set of size at least \(R_q(m;c)\) has an \(m\)-element subset whose
\(q\)-subsets all receive the same color. Thus every red--blue coloring of the
pairs among six objects contains a monochromatic triangle. Van der Waerden's
theorem similarly forces a monochromatic arithmetic progression of any
prescribed finite length in every sufficiently long finite coloring
\cite{Ramsey-GRS-90,GS-90}. These thresholds may be enormous, but they are
finite and computable, unlike a Busy-Beaver horizon.

The guarantee depends on the coding and is weaker than a trading signal. In a
discretized return record, equal symbols at equally spaced dates need not be
consecutive, form a price trend, have a favorable sign, or persist beyond the
sample. Ramsey existence therefore asserts neither statistical dependence nor
causation, predictability, or profit. Nor does it conflict with algorithmic
randomness. A fair-coin Martin-L\"of-random sequence is Borel normal, so every
finite binary word, including arbitrarily long runs and finite chart motifs,
occurs with its fair-coin limiting frequency. Randomness here means escaping
every effective null test and lacking a global effective compression or
successful test-capital process, not local patternlessness
\cite{calude:02,Nies2009}.

Calude and Longo call regularities forced by data volume \emph{Ramsey-type
correlations}. Their incompressibility argument emphasizes that almost all
sufficiently long finite strings resist compression by any fixed positive
fraction while still containing the finite regularities forced by Ramsey
theory~\cite{Calude2016}. This is an epistemic warning, not a theorem that every
sample correlation is false: a Ramsey configuration is not itself a
correlation coefficient, and a pattern observed in one finite history is
compatible with both random and computable continuations.

Trading research compounds this pattern abundance with selection. If all
\(M\) null hypotheses are true, their rejection events are independent, and
each has rejection probability exactly \(\alpha\), the probability of at least
one false rejection is
\(1-(1-\alpha)^M\); if their probabilities are at most \(\alpha\), this is an
upper bound. Without independence, the union bound gives
\(\min\{1,M\alpha\}\) as an upper bound. This
probabilistic multiple-testing problem is distinct from Ramsey's deterministic
existence theorem, but its practical lesson is aligned: a trend or cross-market
relation retained after searching many assets, windows, transformations, lags,
and parameters is a candidate hypothesis, not yet an edge. A defensible FAPP
claim must report the search universe and survive selection correction, costs,
a genuinely untouched holdout or independent replication, and testing outside
the selected regime~\cite{BaileyLopez2014}.

\subsection{What Rice's Theorem Does and Does Not Say}

Rice's theorem makes every nontrivial extensional property of arbitrary partial
computable functions undecidable~\cite{Rice-1953}; in particular, totality is
undecidable in general~\cite{turing-36}. It does not make a fixed total
strategy's gain on a completed finite computable path undecidable, because that
program can simply be simulated. Theorems~\ref{thm:profit-undecidable} and
\ref{thm:market-undecidable} instead give the required promise-domain results
by explicit Halting reductions: the unbounded questions of eventual profit and
future-event occurrence are undecidable even though every code produced by the
reductions is total.

%%------------------------------------------------------------------------------
\section{The No-Arbitrage Foundation}
\label{sec:noarb}
%%------------------------------------------------------------------------------

Here the market is no longer an opponent choosing a response after seeing the
trader's action. Instead, financial admissibility and no-arbitrage delimit
which trader--market pairs are coherent.

\subsection{Discounted Gains and Admissibility}

Let the market be defined on a filtered probability space
\((\Omega,\mathcal{F},\{\mathcal{F}_t\}_{t\in[0,T]},\mathbb{P})\).
Let \(B_t>0\) be a numeraire and \(S_t\) the vector of traded asset prices.
Write \(\overline S_t=S_t/B_t\) for discounted prices. A predictable,
\(\overline S\)-integrable strategy \(H\) has discounted self-financing gain
\begin{equation}
  G_t=(H\mathbin{\cdot}\overline S)_t,\qquad G_0=0.
  \label{eq:gain}
\end{equation}
It is \emph{admissible} if there is a finite constant \(a\) such that
\begin{equation}
  G_t\geq -a\quad\text{for all }t\in[0,T],\quad\mathbb{P}\text{-a.s.}
  \label{eq:admissible}
\end{equation}
This uniform lower bound rules out doubling systems financed by unbounded
credit.

\begin{observation}[Universality Implies Classical Arbitrage]
  \label{thm:arbitrage}
  If an admissible self-financing strategy \(H\) with zero initial capital
  satisfies
  \[
    G_T\geq0\quad\mathbb{P}\text{-a.s.},
    \qquad
    \mathbb{P}(G_T>0)>0,
  \]
  then \(H\) is a classical arbitrage. In particular, a pointwise universal
  strategy is a classical arbitrage whenever its gain is defined on an
  admissible path class carrying \(\mathbb{P}\).
\end{observation}

\begin{proof}
The displayed payoff conditions, together with zero initial cost,
self-financing, and admissibility, are the standard classical-arbitrage
conditions. Pointwise strict positivity on every admissible path implies them,
with \(\mathbb{P}(G_T>0)=1\).
\end{proof}

\begin{remark}[Terminology and Strength]
The pathwise condition is sometimes called a sure or strong arbitrage, but it
is not synonymous with \emph{arbitrage of the first kind}. The latter is the
condition dual to No Unbounded Profit with Bounded Risk (NUPBR/NA1) and is a
distinct, weaker viability concept~\cite{Kardaras2012}. Classical
no-arbitrage alone already excludes the universal strategy in
Observation~\ref{thm:arbitrage}; NFLVR is invoked below because it supplies the
standard martingale-measure formulation.
\end{remark}

\subsection{NA, NUPBR, and Benchmark-Relative Gains}

The relevant no-arbitrage notions do not form the simple chain sometimes
suggested in informal discussions. In the standard semimartingale setting,
NFLVR combines classical NA with NUPBR, whereas NA and NUPBR are in general
logically distinct. NUPBR is equivalent, under standard portfolio hypotheses,
to the existence of a numeraire portfolio whose positive relative-wealth
processes are supermartingales~\cite{KaratzasKardaras2007,Kardaras2012}.
It may hold even when the stronger NFLVR condition and an ELMM fail.

This distinction prevents two overstatements. First, Observation~\ref{thm:arbitrage}
uses only classical NA; NUPBR alone is not silently substituted for it. If a
zero-cost gain were nonnegative at all intermediate times and positive at
terminal time, scaling would violate NUPBR, but terminal pathwise positivity
together with only a strategy-specific lower bound does not by itself supply
that scaling argument. Second, benchmark and stochastic-portfolio approaches
can generate gains \emph{relative} to a numeraire or market portfolio without
producing zero-cost positive gain on every path. Relative arbitrage in diverse
equity markets and real-world benchmark pricing therefore do not contradict
the theorem~\cite{FernholzKaratzasKardaras2005,PlatenHeath2010}.

\subsection{The Martingale Constraint}

For a locally bounded semimartingale price process, the
Delbaen--Schachermayer form of the First Fundamental Theorem states that No
Free Lunch with Vanishing Risk (NFLVR) is equivalent to the existence of an
equivalent local martingale measure (ELMM)
\(\mathbb{Q}\sim\mathbb{P}\) under which \(\overline S\) is a local
martingale~\cite{Harrison1981,Delbaen1994}. For general unbounded
semimartingales the corresponding formulation uses an equivalent
\(\sigma\)-martingale measure~\cite{DelbaenSchachermayer1998}; the corollary below
retains the locally bounded setting.

\begin{corollary}[No Universal Zero-Cost Strategy]
  \label{cor:noarb}
  If an ELMM \(\mathbb{Q}\) exists, no admissible self-financing strategy with
  zero initial capital can satisfy
  \(G_T\geq0\), \(\mathbb{P}\)-almost surely, and
  \(\mathbb{P}(G_T>0)>0\). Consequently, no pointwise universal strategy
  exists among admissible strategies in this stochastic model.
\end{corollary}

\begin{proof}
Under \(\mathbb{Q}\), the admissible stochastic integral
\(G=H\mathbin{\cdot}\overline S\) is a local martingale bounded from below and
therefore a supermartingale. Hence
\(\mathbb{E}_{\mathbb{Q}}[G_T]\leq G_0=0\). If the displayed arbitrage
conditions held, equivalence would give \(G_T\geq0\),
\(\mathbb{Q}\)-almost surely, with
\(\mathbb{Q}(G_T>0)>0\). Consequently
\(\mathbb{E}_{\mathbb{Q}}[G_T]>0\), a contradiction.
\end{proof}

The ELMM conclusion concerns discounted gains under the pricing measure
\(\mathbb{Q}\). It does not say that every risky strategy has non-positive
expected return under the physical measure \(\mathbb{P}\); positive physical
expected returns may compensate risk.

\subsection{Win Rate Is Not Expected Value}

The probability of closing a trade at a profit can greatly exceed \(1/2\)
even in a fair market. Optional stopping gives the precise distinction.

\begin{proposition}[Asymmetric Barriers in a Fair Random Walk]
  \label{prop:optional}
  Let \(X_n\) be a simple symmetric random walk with \(X_0=0\), and for positive
  integers \(a,b\) let
  \[
    \tau=\inf\{n\geq0:X_n=a\text{ or }X_n=-b\}.
  \]
  Then
  \begin{equation}
    \mathbb{P}(X_\tau=a)=\frac{b}{a+b},
    \qquad
    \mathbb{E}[X_\tau]=0.
    \label{eq:gambler}
  \end{equation}
\end{proposition}

\begin{proof}
For this finite interval, \(\tau<\infty\) almost surely and
\(|X_{n\wedge\tau}|\leq\max\{a,b\}\). Optional stopping gives
\(\mathbb E[X_{n\wedge\tau}]=0\), and bounded convergence yields
\(\mathbb E[X_\tau]=0\)~\cite{Doob1953}. If
\(p=\mathbb{P}(X_\tau=a)\), then
\(0=\mathbb{E}[X_\tau]=ap-b(1-p)\), which gives
\(p=b/(a+b)\).
\end{proof}

With a take-profit of \(a=1\) and a stop-loss of \(b=9\), the win probability
is \(0.9\), but the expected payoff is zero: nine frequent unit gains are
offset by one nine-unit loss in expectation. Thus prediction accuracy, trade
win rate, and expected discounted profit are three different quantities.

%%------------------------------------------------------------------------------
\section{The Combinatorial Perspective}
\label{sec:nfl}
%%------------------------------------------------------------------------------

Here the opposing side is an ensemble of market records rather than an
adaptive market. The optimization NFL theorem and the elementary prediction
result needed for market-direction claims should be distinguished.

\begin{theorem}[Classical Optimization NFL, Wolpert--Macready]
  \label{thm:nfl}
  Let \(X\) and \(Y\) be finite and give the function class \(Y^X\) its uniform
  law. For any two non-revisiting black-box algorithms and any fixed number
  \(q\leq |X|\) of queries, the resulting \(q\)-term sequences of observed
  values have the same distribution. Consequently, every performance
  functional depending only on that observed value sequence has the same
  average for both algorithms
  \cite{Wolpert:1997:NFL:2221336.2221408}.
\end{theorem}

The theorem concerns oracle-based optimization over function spaces. It does
not, by itself, imply a statement about self-financing gains on stochastic
price paths. A sharpening makes its structural assumption explicit: for a
uniform distribution on a finite function class, the NFL equality for all
non-revisiting algorithms holds precisely when the class is closed under
permutations of the search domain~\cite{SchumacherVoseWhitley2001}. Financial
path classes are generally not permutation-closed: temporal order, continuity,
financing, and no-arbitrage restrictions distinguish one permutation from
another. Thus a financial edge, when present, resides in exactly such
symmetry-breaking structure and the inductive bias used to model it.

For binary directional prediction the relevant result is simpler and can be
derived directly.

\begin{proposition}[Uniform Binary Prediction]
  \label{prop:binary}
  Let \(X_1,\ldots,X_T\) be uniformly distributed on \(\{0,1\}^T\). A causal
  deterministic predictor chooses
  \(\widehat X_t=f_t(X_1,\ldots,X_{t-1})\). Its mean accuracy
  \[
    A_T=\frac1T\sum_{t=1}^T
      \mathbf{1}\{\widehat X_t=X_t\}
  \]
  satisfies \(\mathbb{E}[A_T]=1/2\). The same holds for a randomized predictor
  driven by a private random seed independent of
  \((X_1,\ldots,X_T)\).
\end{proposition}

\begin{proof}
Conditional on every history \(X_1,\ldots,X_{t-1}\), the next bit \(X_t\) is
an independent fair bit. Therefore
\(\mathbb{P}(\widehat X_t=X_t\mid X_1,\ldots,X_{t-1})=1/2\).
Linearity of expectation gives the result. In the randomized case, conditioning
additionally on the independent private seed leaves \(X_t\) fair given the
history and proves the same equality.
\end{proof}

Consequently, no predictor can achieve expected directional accuracy greater
than \(1/2\) under the uniform distribution over all binary sequences; in
particular, none can exceed \(1/2\) on every sequence. This is an NFL symmetry
statement, not a model of actual markets. Financial sign sequences are not
uniform in general, and a directional forecast does not determine trading
profit because payoff sizes, financing, and stopping rules matter
(Proposition~\ref{prop:optional}).

The defensible practical conclusion is conditional: a claim of repeatable
success must identify non-uniform structure in the physical measure
\(\mathbb{P}\), such as a trend, mean reversion, a risk premium, information,
or an institutional constraint that breaks the NFL symmetry. The hypothesized
edge concerns a restricted environment, not universal algorithmic
superiority.

%%------------------------------------------------------------------------------
\section{Practical Robustness Diagnostics: Time Reversal and Permutations}
\label{sec:criterion}
%%------------------------------------------------------------------------------

This section turns the trader--market leitmotif into a falsification tool. Hold
the trader fixed and transform only the market record. Among the countable
family of computably specified stress transformations, time reversal is one of
the simplest, particularly for detecting dependence on directional
persistence. It is not a further impossibility theorem on the level of the
Halting reductions or no-arbitrage.

Let \(\mathcal{M}\) be a space of strictly positive discounted price paths
\(P\colon[0,T]\to\mathbb{R}_{>0}^n\). The time reversal of \(P\) is
\begin{equation}
  \widetilde{P}(t)=P(T-t), \qquad t\in[0,T].
  \label{eq:timereversal}
\end{equation}
Literal reversal reuses the observed price levels and preserves the magnitudes
of log-price moves while reversing their order and signs. It is not the same
as merely negating arithmetic returns: if a forward gross return is \(R\), the
corresponding reversed gross return is \(1/R\). In an empirical implementation
the strategy must be rerun causally from scratch on the reversed record;
reversing the original orders or realized profit would leak future information.
If the initial level is prescribed, literal reversal can leave the market
class. One may instead test componentwise normalized reversal,
\[
  P_i^{\mathrm{rev}}(t)=P_i(0)\frac{P_i(T-t)}{P_i(T)},
\]
provided it remains in the claimed class. Scaling can change the actions of a
strategy that is not scale invariant; it is therefore a distinct test.
For a causal strategy \(\sigma\), let \(\Pi_t[\sigma;P]\) denote its discounted
gain from zero initial capital through time \(t\) on \(P\). Call \(\sigma\)
\emph{pathwise admissible} on \(\mathcal M\) if some finite \(C\) satisfies
\(\Pi_t[\sigma;P]\geq-C\) for every \(P\in\mathcal M\) and
\(t\in[0,T]\).

\begin{definition}[Pathwise Universal Strategy]
\label{def:universal}
  A causal strategy \(\sigma\) is \emph{pathwise universal} on
  \(\mathcal{M}\) if it is self-financing and pathwise admissible there and
  satisfies \(\Pi_T[\sigma;P]>0\) for every \(P\in\mathcal{M}\).
\end{definition}

\begin{criterion}[Reversal Falsification Test]
  \label{crit:tr}
  Suppose \(\mathcal{M}\) is closed under time reversal. If there exists
  \(P\in\mathcal{M}\) such that
  \begin{equation}
    \min\{\Pi_T[\sigma;P],\Pi_T[\sigma;\widetilde P]\}\leq 0,
    \label{eq:reversaltest}
  \end{equation}
  then \(\sigma\) is not pathwise universal on \(\mathcal{M}\).
\end{criterion}

\begin{proof}
Both \(P\) and \(\widetilde P\) belong to \(\mathcal{M}\). A non-positive gain
on either member of the pair contradicts Definition~\ref{def:universal}.
\end{proof}

One non-winning path refutes a claim quantified over every path. Reversal
supplies a paired candidate, not a proof that one member must fail: a causal
strategy may take different positions on the reversed history.
Diagonalization and no-arbitrage provide the genuine impossibility results
under their respective assumptions.

\paragraph{Permuting a finite record}

Time reversal is a special permutation of a sampled record. For a finite
record of discounted price levels \(P=(P_0,\ldots,P_N)\), let \(\pi\) be
a permutation of \(\{0,\ldots,N\}\) and define
\begin{equation}
  (P^\pi)_t=P_{\pi(t)},\qquad t=0,\ldots,N.
  \label{eq:price-permutation}
\end{equation}
Thus \(\pi(t)=N-t\) gives reversal. A fixed algorithm or a random shuffle
can choose \(\pi\). For vector-valued prices, each complete price vector
moves together. Let \(\mathcal M_N\) be the claimed class of such discrete
records and \(\Pi_N[\sigma;P]\) the discounted gain from zero initial capital
when the same causal rule is rerun from its prescribed initial state on \(P\).
Use the discrete version of Definition~\ref{def:universal} on this class.

\begin{criterion}[Permutation Falsification Test]
  \label{crit:permutation}
  A pathwise universal strategy on \(\mathcal M_N\) must satisfy
  \begin{equation}
    \begin{aligned}
      &\Pi_N[\sigma;P^\pi]>0\\
      &\quad\text{whenever }P\in\mathcal M_N\text{ and }P^\pi\in\mathcal M_N.
    \end{aligned}
    \label{eq:permutationtest}
  \end{equation}
  One permitted permutation with zero or negative gain refutes that universal
  claim. The criterion requires positivity, not equality of profits across
  permutations.
\end{criterion}

\begin{proof}
Every retained \(P^\pi\) is a member of the claimed market class, so strict
profit on every member requires strict profit on this one.
\end{proof}

The shuffle rearranges price levels, not returns or previously chosen orders.
At each date the rerun uses only the history available in the new order.
Permuting levels preserves their multiset but generally changes the returns
and their magnitudes. It can also violate a fixed initial-price condition,
a bound on price moves, or another restriction defining \(\mathcal M_N\).
Only transformations that remain within that class test its universal claim;
closure under all permutations is unnecessary. Failure after destroying a
trend outside a strategy's stated domain does not refute a claim confined to
that domain. Passing finitely many shuffles of one record does not establish
universal profitability on a larger path class.

Unlike the permutation closure used in optimization NFL
(Section~\ref{sec:nfl}), this diagnostic assumes no uniform distribution or
equality of average performance. Randomness in the choice of stress records
alone supplies neither a market law nor a statistical significance level.

Greg Egan's \emph{Permutation City} (1994)~\cite{permutationcity} provides a
literary analogy by exploring the temporal ordering of simulated computational
states. Here permutations change the trader's observation sequence; the causal
rule is rerun on each record, and Criterion~\ref{crit:permutation} follows from
the universal-profit requirement.

\begin{remark}[Computable Stress Tests Beyond Reversal]
\label{rem:stressfamily}
Time reversal is one of countably many computable probes: sign reversal,
regime splicing, jumps, delays, volatility rescaling, financing, and transaction
costs. Since program totality is undecidable, an operational battery must select
known-terminating transformations whose outputs remain in the claimed market
class. Here ``test'' means empirical stress, not a Martin-L\"of test: a uniformly
computably enumerable sequence of open sets with effective measure bounds.
Reversal defines no null set and certifies no randomness; passing it shows only
tested robustness, not universality.
\end{remark}

\begin{remark}[Filtrations and Time Reversal]
Stochastic model classes are not automatically reversal-closed:
\(P(T-t)\) is generally not adapted to the forward filtration, and reversing a
diffusion requires additional hypotheses~\cite{HaussmannPardoux1986}. For
example, if
\[
  X_t=\log(S_t/S_0)=mt+\sigma W_t,
\]
then \(X_{T-t}-X_T\), under its reversed filtration, has drift \(-m\) and
volatility \(\sigma\), so a class containing both drift signs is closed in law
under reversed increments. Physically, \(m=\mu-\sigma^2/2\); under a
risk-neutral measure the discounted stock may be a martingale while its
logarithm retains It\^o drift \(-\sigma^2/2\). Neither an ELMM nor a driftless
arithmetic-Brownian special case establishes detailed balance in general. The
diagnostic remains pathwise; Section~\ref{sec:noarb} gives the probabilistic
constraints.
\end{remark}

%%------------------------------------------------------------------------------
\section{Case Study: The Wheel Options Strategy}
\label{sec:wheel}
%%------------------------------------------------------------------------------

The Wheel makes the trader--market configuration concrete. It sells a
cash-secured put and, after assignment, sells a covered call against the
acquired stock. For one matched pair of European options, with strike \(K\)
and maturity \(T\), the terminal-payoff identity is
\begin{equation}
  S_T-(S_T-K)^+
  =K-(K-S_T)^+
  =\min\{S_T,K\}.
  \label{eq:putcall}
\end{equation}
The cash leg pays \(K\) at maturity and must be financed at its present value;
dividends and other cash flows must be matched separately. This is a local
payoff comparison, not a model of the rolling Wheel. American exercise,
assignment, changing strikes and maturities, roll rules, spreads, taxes, and
financing make the realized strategy path dependent. Each short-option leg
nevertheless carries the associated downside exposure.

Option selling can nevertheless have positive mean return under a physical
law. The
difference between risk-neutral and physical expectations of future variance
is commonly called the variance risk premium, with sign conventions varying
across the literature, and option-market evidence documents compensation for
bearing volatility and jump risk
\cite{BakshiKapadia2003,CarrWu2009}. Such compensation is empirical and
state-dependent: its sign and magnitude depend on option prices, strikes,
transaction costs, and the physical distribution.

The local payoff comparison suggests four stress scenarios for the rolling strategy. A crash
after assignment produces essentially one-for-one downside below \(K\). A
persistent decline can make repeated premiums too small to offset the renewed
short-put exposure. In a large rally, assignment and call exercise can still
leave a positive absolute profit, but the cap at \(K\) can underperform a
buy-and-hold benchmark. Finally, capital exhaustion is absorbing for a
self-financing trader: if \(W_{t+1}=W_t(1+fR_{t+1})\) and
\(1+fR_{t+1}=0\), then \(W_{t+1}=0\) unless external capital is injected. A
fully cash-secured, unlevered implementation has bounded loss, but an
underlying that falls to zero can consume almost all committed capital net of
premium.

A rally/crash reversal is therefore a useful paired stress test, although the
payoff identity, rather than reversal alone, supplies the economic mechanism.
Repeated small premiums also need not imply satisfactory time-average growth:
\[
  \frac1T\log\frac{W_T}{W_0}
  =\frac1T\sum_{t=1}^T\log(1+fR_t),
\]
so frequent gains can coexist with poor long-run growth when rare losses are
large~\cite{Peters2019}. The Wheel may be a defensible FAPP strategy under a
specified risk premium, capital policy, cost model, and benchmark; it is not a
pathwise profit guarantee.

%%------------------------------------------------------------------------------
\section{Universal Portfolios and Distance from Universality}
\label{sec:cover}
%%------------------------------------------------------------------------------

Here the trader changes the claim made against the market side: it seeks small
regret to a benchmark portfolio class, not positive profit on every path. The
word ``universal'' also appears in a celebrated positive theorem.
For price-relative vectors \(x_t\in\mathbb{R}_{>0}^m\), a constant-rebalanced
portfolio \(b\) has wealth
\[
  S_T(b)=\prod_{t=1}^T b^\mathsf{T}x_t.
\]
Cover's universal portfolio is nonanticipating and, under the conditions of
his theorem, has the same asymptotic exponential growth rate as the best
constant-rebalanced portfolio chosen in hindsight on every individual bounded
sequence~\cite{Cover1991}. In regret notation,
\begin{equation}
  \frac1T\left[
  \log\max_{b}S_T(b)-\log S_T^{\mathrm{UP}}
  \right]\longrightarrow0.
  \label{eq:cover}
\end{equation}

Equation~\eqref{eq:cover} is a relative guarantee, not an absolute-profit
guarantee. If every constant-rebalanced portfolio loses money on a path, the
universal portfolio may lose as well while retaining vanishing per-period
regret. Cover's result therefore does not contradict
Corollary~\ref{cor:noarb}; it illustrates a useful benchmark-relative meaning
of ``universal'' once the comparison class is specified.

Online regret methods also yield financial guarantees of another kind.
DeMarzo, Kremer, and Mansour show that gradient strategies designed to
minimize asymptotic regret imply trading strategies that provide
arbitrage-based option price bounds without assuming continuous price paths,
complete markets, or a pricing kernel; the bounds depend on realized quadratic
variation
\cite{DeMarzoKremerMansour2016}. This is robust pricing and hedging, not
guaranteed positive profit or next-move prediction.

\paragraph{A Hierarchy of Attainable Claims}

The useful alternatives to pathwise universality are not weaker points on one
single order; they change the benchmark, horizon, or probability quantifier.
Table~\ref{tab:hierarchy} records the distinctions.

\begin{table*}[ht]
\caption{\label{tab:hierarchy}Different meanings of a successful or universal strategy.}
\centering
\begin{ruledtabular}
\begin{tabular}{p{0.22\textwidth}p{0.36\textwidth}p{0.34\textwidth}}
\textbf{Notion} & \textbf{Guarantee} & \textbf{Status} \\
\colrule
Pathwise universal profit
& Positive discounted zero-cost gain on every admissible path
& Excluded among admissible strategies in the specified stochastic model under classical NA \\
Relative arbitrage
& Outperformance of a specified market or numeraire portfolio
& Possible under structural market assumptions \\
Cover universality
& Vanishing per-period log regret to the best constant-rebalanced portfolio
& Pathwise but benchmark-relative \\
FAPP or statistical edge
& Positive performance under a specified physical law, data protocol, and horizon
& Conditional and empirically falsifiable \\
\end{tabular}
\end{ruledtabular}
\end{table*}

Superhedging duality gives another quantitative notion: once a contingent
claim and an admissible model class are fixed, its superhedging price is the
least initial capital needed for a pathwise guarantee. Requiring zero initial
capital while demanding a strictly positive payoff is precisely the
degenerate, arbitrage-producing endpoint ruled out above.

%%------------------------------------------------------------------------------
\section{FAPP Strategies and Their Honest Scope}
\label{sec:fapp}
%%------------------------------------------------------------------------------

Although pathwise universal profit is excluded, a strategy may be profitable
\emph{for all practical purposes} (FAPP, following Bell~\cite{bell-a}) relative
to a physical law \(\mathbb P\), a benchmark, and a finite horizon. Market
making may earn spreads when adverse selection and inventory risk are
controlled, and Kelly investing maximizes expected logarithmic growth when the
relevant distribution is sufficiently well specified~\cite{Kelly1956}. Neither
is an unconditional guarantee.

Here ``insight'' means a falsifiable restriction on the market side. In actual
trading it usually appears as a heuristic: persistence, mean reversion,
valuation correction, or compensation for an identified risk. The fact that a
heuristic is algorithmic does not disqualify it; its edge is a regime-dependent
empirical hypothesis. A repeatable ex ante claim must therefore specify its
market class, horizon, costs, benchmark, and validation protocol.

Only lawfully acquired and used information is included in this practical
strategy class; an informational advantage remains short of a pathwise
guarantee.

The unit of performance matters as well. For a funded portfolio, let
\(R^{\mathrm{nom}}\) be nominal total return and \(\pi\) inflation over the same
interval under a stated price index. Then
\begin{equation}
  1+R^{\mathrm{real}}
  =\frac{1+R^{\mathrm{nom}}}{1+\pi}.
  \label{eq:real-return}
\end{equation}
A positive nominal return can therefore be a purchasing-power loss. The price
index is an economic deflator, not necessarily a tradable numeraire for the
FTAP; the consumption benchmark and financing account serve different
purposes.

Productivity growth may support positive real returns when its benefits reach
investors as cash flows~\cite{Solow1957}. Prices, discount rates, and the
displacement of incumbents intervene~\cite{Schumpeter1942}; aggregate growth
therefore establishes neither next-move prediction nor active alpha.

Equilibrium theory likewise permits only conditional edges. Grossman and
Stiglitz show that perfectly informationally efficient markets are inconsistent
with costly information acquisition: if prices revealed all information
without reward, agents would not pay to produce it~\cite{GrossmanStiglitz1980}.
Conversely, under its common-prior, common-knowledge, risk-aversion, and
Pareto-optimality assumptions, the Milgrom--Stokey no-trade theorem limits trade
based on private information alone~\cite{MilgromStokey1982}. A complementary
market-selection result is also conditional: among agents in Sandroni's
complete-market model who share an intertemporal discount factor and choose
saving endogenously, accurate forecasters prosper and inaccurate ones are driven
out~\cite{Sandroni2000}. These equilibrium results depend on their model
assumptions; none is a pathwise trading guarantee or a substitute for the FTAP.

Deployment can also change the market side. A large, crowded, or public
strategy consumes liquidity, moves quotes, attracts imitation or opposition,
and changes the distribution on which it was estimated. Lo's Adaptive Markets
Hypothesis provides this evolutionary interpretation~\cite{Lo2004}: competition
and impact may erode an edge without producing the tailored countermarket of
Theorem~\ref{thm:counterpath}. Feedback invalidates scale-free extrapolation; it
does not imply that every large strategy loses.

Finally, the finite-pattern trap of Section~\ref{sec:ramseytrap} becomes
operational in strategy discovery. Trying many specifications and retaining the
best backtest turns its Sharpe ratio into an order statistic. The Deflated
Sharpe Ratio corrects for selection, non-normality, and the number of
trials~\cite{BaileyLopez2014}; untouched validation, costs, and deployment
analysis remain necessary. These controls do not prove an edge, but make a
finite-sample FAPP claim more honestly falsifiable.

%%------------------------------------------------------------------------------
\section{Conclusion}
\label{sec:conclusion}
%%------------------------------------------------------------------------------

Everything returns to the same pair: a trader and a market. For every fixed
total deterministic computable trader, Theorem~\ref{thm:counterpath} constructs
a fixed positive computable price path whose next bounded proportional move
opposes every nonzero position and remains nonzero during inactivity.
The quantifiers are
\(\forall\sigma\,\exists\mathcal E_\sigma\): no physical market must observe
the deployed orders, the always-flat trader earns zero, and every nonzero
position loses on its trader-relative countermarket. This direct diagonal
result uses no Halting-Problem reduction and identifies no single path that
defeats every trader.

The remaining limits concern different market sides. Gold precludes identifying
every computable binary generator in the limit, while next-bit diagonalization
defeats each fixed computable predictor on its predictor-relative record.
Halting reductions preclude uniform certification of eventual profit or encoded
future events, and Busy Beaver excludes a computable description-size waiting
horizon despite finitely solvable low-complexity cases. Under a stated
computable reference measure, Martin-L\"of randomness blocks unbounded success
by effective test-capital processes. Ramsey regularity explains why finite
structured subsequences nevertheless remain abundant, including in random
records. No-arbitrage and uniform binary averaging supply still different
boundaries. Their conclusions are complementary because their inputs,
quantifiers, and success criteria differ.

Positive results arise only after restricting the claim. The Wheel may earn a
risk premium while retaining put-like downside; Cover guarantees asymptotic
performance relative to a benchmark class; and FAPP strategies may exploit a
specified regime, lawful information, or compensated risk. Such performance
must be evaluated after costs and financing, against an explicit benchmark and,
where relevant, in purchasing-power terms. Productivity-supported real growth
is one possible regime hypothesis, not a next-move predictor or a guarantee of
active alpha.

Any repeatable-success claim must therefore state its market restriction,
probability law, benchmark, admissibility and financing conditions, horizon,
data protocol, and deployment effects. Time reversal and other within-class
permutations are robustness probes: one nonpositive gain refutes an all-path
claim, while passing a finite set of probes does not certify universality.
Useful strategies may exist, but no total deterministic computable
self-financing rule guarantees strict gain on every positive computable price
path; under an arbitrage-free stochastic model, no admissible zero-cost rule
delivers nonnegative discounted gain almost surely and positive gain with
positive probability.

\begin{acknowledgments}
This research was funded in whole or in part by the Austrian Science Fund
(FWF) [Grant DOI: 10.55776/PIN5424624]. For open access purposes, the author
has applied a CC BY public copyright license to any author accepted manuscript
version arising from this submission. The author acknowledges TU Wien
Bibliothek for financial support through its Open Access Funding Programme.

OpenAI Codex was used to assist with literature organization, LaTeX
editing, and checks of derivations. The author specified the scientific
direction and prompts; model output retained in the manuscript was checked
against the cited sources and explicit calculations. The author accepts full
responsibility for the final text.
\end{acknowledgments}

\section*{Data Availability}
This is a purely mathematical work, and no data or software were created or
analyzed in this study.

%%------------------------------------------------------------------------------
\clearpage
\bibliography{svozil}
\end{document}
